% 康奈尔听课笔记模板（配色取自王道课件）
\documentclass[10.5pt,a4paper,oneside]{article}
\usepackage[a4paper,top=2.3cm,bottom=2.2cm,left=5.7cm,right=1.5cm,
  marginparwidth=4.25cm,marginparsep=0.55cm,headheight=14pt,footskip=1.1cm]{geometry}
\usepackage{xeCJK}
\usepackage{xeCJKfntef}
\xeCJKDeclareCharClass{FullLeft}{"201C,"2018}
\xeCJKDeclareCharClass{FullRight}{"201D,"2019}
\xeCJKDeclareCharClass{CJK}{"2460->"24FF,"2605,"2606}
\setmainfont{Noto Sans CJK SC}[BoldFont={Noto Sans CJK SC Bold},ItalicFont={Noto Sans CJK SC},ItalicFeatures={FakeSlant=0.18}]
\setsansfont{Noto Sans CJK SC}[BoldFont={Noto Sans CJK SC Bold}]
\setmonofont{DejaVu Sans Mono}[Scale=0.9]
\setCJKmainfont{Noto Sans CJK SC}[BoldFont={Noto Sans CJK SC Bold},ItalicFont={Noto Sans CJK SC},ItalicFeatures={FakeSlant=0.18}]
\setCJKsansfont{Noto Sans CJK SC}[BoldFont={Noto Sans CJK SC Bold}]
\setCJKmonofont{Noto Sans Mono CJK SC}[Scale=0.9]
\linespread{1.32}
\usepackage{amsmath,amssymb}
\usepackage[dvipsnames,table]{xcolor}
\usepackage{graphicx}
\usepackage{adjustbox}
\usepackage{tikz}
\usetikzlibrary{calc}
\usepackage[most]{tcolorbox}
\usepackage{eso-pic}
\usepackage{fancyhdr}
\usepackage{titlesec}
\usepackage{needspace}
\usepackage{enumitem}
\usepackage{xltabular,booktabs,array,colortbl}
\usepackage{marginnote}
\reversemarginpar
\renewcommand*{\raggedleftmarginnote}{}
\renewcommand*{\raggedrightmarginnote}{}
\usepackage{caption}
\usepackage{titletoc}
\usepackage{fancyvrb}
\usepackage{framed}
\usepackage[section]{placeins}
\usepackage{hyperref}
\hypersetup{colorlinks=true,linkcolor=wdblue,urlcolor=wdblue,pdftitle={3.2
课后习题：队列},pdfauthor={听课笔记}}

% ---------- 课件配色 ----------
\definecolor{wdblue}{HTML}{4472C4}
\definecolor{wdlightblue}{HTML}{5B9BD5}
\definecolor{wdtitle}{HTML}{4A86C8}
\definecolor{wddark}{HTML}{1F3864}
\definecolor{wdorange}{HTML}{ED7D31}
\definecolor{wdgray}{HTML}{A5A5A5}
\definecolor{wdgold}{HTML}{FFC000}
\definecolor{wdgreen}{HTML}{70AD47}
\definecolor{wddeepgreen}{HTML}{3F7F2A}
\definecolor{wdred}{HTML}{C00000}
\definecolor{wdpurple}{HTML}{7030A0}
\definecolor{cuebg}{HTML}{F2F6FC}
\definecolor{codebg}{HTML}{F7F9FC}
\color[HTML]{262626}

% ---------- 行内样式 ----------
\newcommand{\hlred}[1]{\textcolor{wdred}{\bfseries #1}}
\newcommand{\hlblue}[1]{\textcolor{wdblue}{\bfseries #1}}
\newcommand{\hlgreen}[1]{\textcolor{wddeepgreen}{\bfseries #1}}
\newcommand{\hlorange}[1]{\textcolor{wdorange}{\bfseries #1}}
\newcommand{\hlpurple}[1]{\textcolor{wdpurple}{\bfseries #1}}
\newcommand{\uln}[1]{\CJKunderline[format=\color{wdred},thickness=0.9pt,depth=0.28em]{#1}}
\newcommand{\wavy}[1]{\CJKunderwave[format=\color{wdorange}]{#1}}
\newcommand{\hlmark}[1]{\CJKunderline[format=\color{wdgold!55},thickness=0.55em,depth=-0.05em]{#1}}
\newcommand{\tagbox}[2]{\tcbox[on line,colback=#1,colframe=#1,coltext=white,boxsep=0.5pt,left=3pt,right=3pt,top=0.5pt,bottom=0.5pt,arc=2pt,fontupper=\small\bfseries]{#2}}

% ---------- 康奈尔版式背景：左侧线索栏 ----------
\newlength{\cuestripw}\setlength{\cuestripw}{5.15cm}
\AddToShipoutPictureBG{%
  \begin{tikzpicture}[remember picture,overlay]
    \fill[cuebg] (current page.north west) rectangle ($(current page.south west)+(\cuestripw,0)$);
    \draw[wdblue,line width=1.2pt] ($(current page.north west)+(\cuestripw,-1.55cm)$) -- ($(current page.south west)+(\cuestripw,1.55cm)$);
    \node[anchor=west,text=wdblue,font=\small\bfseries] at ($(current page.north west)+(0.75cm,-1.25cm)$) {线索 · 问题};
    \node[anchor=west,text=wdblue,font=\small\bfseries] at ($(current page.north west)+(\cuestripw+0.55cm,-1.25cm)$) {笔记 · 老师主线 + 补充讲解};
  \end{tikzpicture}}

% ---------- 页眉页脚 ----------
\pagestyle{fancy}\fancyhf{}
\renewcommand{\headrulewidth}{0pt}
\fancyhead[R]{\small\color{wdgray}第3章 栈、队列和数组}
\fancyfoot[C]{\small\color{wdgray}\thepage}
\fancyfoot[R]{\small\color{wdorange}3.2 课后习题：队列 · 听课笔记}
\fancyfoot[L]{\small\hyperlink{toc}{\textcolor{wdblue}{\ensuremath{\uparrow}\,回到目录}}}
\fancypagestyle{plain}{\fancyhf{}\fancyfoot[C]{\small\color{wdgray}\thepage}\fancyfoot[R]{\small\color{wdorange}3.2
课后习题：队列 · 听课笔记}}

% ---------- 标题：仿课件的蓝色标题 + 六色条 ----------
\newcommand{\sixbar}{\par\nointerlineskip\vspace{3pt}\centering
  \begin{tikzpicture}\foreach \c [count=\i] in {wdlightblue,wdorange,wdgray,wdgold,wdblue,wdgreen}
     {\fill[\c] ({(\i-1)*1.05},0) rectangle ({(\i-1)*1.05+0.92},0.11);}\end{tikzpicture}\par}
\titleformat{\section}[block]{\centering\Large\bfseries\color{wdtitle}}{}{0pt}{}[\sixbar]
\titlespacing*{\section}{0pt}{16pt}{10pt}
\pretocmd{\section}{\Needspace{8\baselineskip}}{}{}
\titleformat{\subsection}{\Needspace{7\baselineskip}\large\bfseries\color{wddark}}{\textcolor{wdblue}{\rule[-0.15em]{4pt}{1.05em}}\hspace{6pt}}{0pt}{}
\titlespacing*{\subsection}{0pt}{12pt}{5pt}
\titleformat{\subsubsection}{\Needspace{5\baselineskip}\normalsize\bfseries\color{wdblue}}{\textcolor{wdorange}{\large\ensuremath{\blacktriangleright}}\hspace{4pt}}{0pt}{}
\titlespacing*{\subsubsection}{0pt}{8pt}{3pt}
\setcounter{secnumdepth}{0}
\setlength{\parindent}{0pt}\setlength{\parskip}{5pt}
\setlist{nosep,leftmargin=1.5em,itemsep=2pt}
\setlist[itemize,1]{label=\textcolor{wdblue}{\small\ensuremath{\bullet}}}
\setlist[itemize,2]{label=\textcolor{wdorange}{\scriptsize\ensuremath{\blacktriangleright}}}
\setlist[enumerate,1]{label=\textcolor{wdblue}{\bfseries\arabic*.}}
\providecommand{\tightlist}{\setlength{\itemsep}{2pt}\setlength{\parskip}{0pt}}

% ---------- 目录 ----------
\setcounter{tocdepth}{1}
\titlecontents{section}[0pt]{\addvspace{5pt}\bfseries\color{wddark}}{}{}{\titlerule*[0.7pc]{\color{wdgray!60}.}\contentspage}
\titlecontents{subsection}[1.6em]{\addvspace{1pt}\small\color{black!75}}{}{}{\titlerule*[0.7pc]{\color{wdgray!40}.}\contentspage}
\newcommand{\tocback}{\par\vspace{-3pt}{\hfill\footnotesize\hyperlink{toc}{\textcolor{wdgray}{\ensuremath{\uparrow}\,回到目录}}}\par\vspace{-4pt}}

% ---------- 线索栏 ----------
\newcommand{\cueblock}[2]{\Needspace{4\baselineskip}\leavevmode#1{#2}\par\vspace{-\baselineskip}\vspace{-\parskip}}
% 线索直接贴在段首行的左侧（不用 marginnote，避免分页时偶尔落进正文栏）
\newcommand{\cuebox}[1]{\leavevmode\llap{\smash{\begin{minipage}[t]{\marginparwidth}\raggedright\small\linespread{1.18}\selectfont #1\end{minipage}}\hspace{\marginparsep}}}
\newcommand{\cuenote}[1]{\cuebox{{\color{wdblue}\bfseries Q\,}\color{wddark}#1}}
\newcommand{\cuekey}[1]{\cuebox{\color{wdred}\bfseries #1}}

% ---------- 讲解框 ----------
\tcbset{wdbox/.style={enhanced,breakable,arc=2pt,boxrule=0pt,leftrule=3pt,left=7pt,right=7pt,top=5pt,bottom=5pt,
   before skip=7pt,after skip=7pt,fonttitle=\bfseries\small,attach title to upper={\quad},coltitle=#1,
   colframe=#1,colback=#1!6}}
\newtcolorbox{defbox}[1][]{enhanced,breakable,arc=3pt,boxrule=0.8pt,colframe=wdblue,colback=wdblue!4,
   left=8pt,right=8pt,top=6pt,bottom=6pt,fonttitle=\bfseries\small,coltitle=white,colbacktitle=wdblue,
   title={定义#1},before skip=8pt,after skip=8pt}
\newtcolorbox{thmbox}[1]{enhanced,breakable,arc=3pt,boxrule=0.8pt,colframe=wdblue,colback=wdblue!4,
   left=8pt,right=8pt,top=6pt,bottom=6pt,fonttitle=\bfseries\small,coltitle=white,colbacktitle=wdblue,
   title={#1},before skip=8pt,after skip=8pt}
\newtcolorbox{bridgebox}{wdbox=wdblue,title={承上启下}}
\newtcolorbox{plainbox}{wdbox=wddeepgreen,colframe=wdgreen,colback=wdgreen!8,title={大白话}}
\newtcolorbox{deepbox}{wdbox=wdpurple,title={背后的思想}}
\newtcolorbox{exambox}{wdbox=wdred,title={考点}}
\newtcolorbox{pitfallbox}{wdbox=wdorange,title={易错提醒}}
\newtcolorbox{bookbox}{wdbox=wddark,colback=wdgray!10,colframe=wdgray,title={王道书补充}}
\newtcolorbox{teacherbox}{wdbox=wdorange,colback=wdgold!8,colframe=wdgold,coltitle=wdorange,title={老师原话}}
\newtcolorbox{fixbox}{wdbox=wdred,colback=wdred!4,title={勘误 / 说明}}
% 习题笔记用
\newtcolorbox{questionbox}[1][]{enhanced,breakable,arc=3pt,boxrule=0.8pt,colframe=wddark,colback=white,
   left=8pt,right=8pt,top=6pt,bottom=6pt,fonttitle=\bfseries\small,coltitle=white,colbacktitle=wddark,
   title={#1},before skip=10pt,after skip=6pt}
\newtcolorbox{videobox}{wdbox=wdorange,colback=wdgold!8,colframe=wdgold,coltitle=wdorange,title={视频讲解}}
\newtcolorbox{jiexibox}{wdbox=wddark,colback=wdgray!10,colframe=wdgray,title={书上解析}}
\newtcolorbox{addbox}{wdbox=wdpurple,title={我的补充}}
\newtcolorbox{recallbox}{wdbox=wdblue,title={回顾正课}}
\newtcolorbox{watchbox}{enhanced,breakable,arc=3pt,boxrule=1pt,colframe=wdorange,colback=wdgold!12,
   fonttitle=\bfseries\small,coltitle=white,colbacktitle=wdorange,title={建议去看视频},left=8pt,right=8pt,before skip=8pt,after skip=8pt}
\newtcolorbox{summarybox}{enhanced,arc=3pt,
   boxrule=1pt,colframe=wdblue,colback=white,colbacktitle=wdblue,coltitle=white,fonttitle=\bfseries,
   title={总结区 · 用自己的话复述这一节},left=10pt,right=10pt,top=8pt,bottom=8pt,before skip=14pt}

% ---------- 代码 ----------

\makeatletter
\@ifundefined{Shaded}{}{\renewenvironment{Shaded}{\def\FrameCommand{\fboxsep=5pt\fboxrule=0.5pt\fcolorbox{wdblue!35}{codebg}}%
   \MakeFramed{\advance\hsize-\width\FrameRestore}\small}{\endMakeFramed}}
\makeatother
\fvset{fontsize=\small}

% ---------- 表格 ----------
\renewcommand{\arraystretch}{1.42}
\newcolumntype{Y}{>{\raggedright\arraybackslash}X}

% ---------- 图片 ----------
% 课件截图做成浮动体：LaTeX 会自动用正文填满页面，避免大块空白
\setcounter{topnumber}{3}\setcounter{bottomnumber}{2}\setcounter{totalnumber}{4}
\renewcommand{\topfraction}{0.85}\renewcommand{\bottomfraction}{0.6}\renewcommand{\textfraction}{0.12}\renewcommand{\floatpagefraction}{0.75}
\setlength{\textfloatsep}{10pt plus 3pt minus 2pt}\setlength{\floatsep}{8pt plus 2pt}
\newcommand{\slidefig}[3]{\begin{figure}[!ht]\centering\begin{tcolorbox}[enhanced,width=#2\linewidth,arc=2pt,boxrule=0.6pt,colframe=wdgray!70,
   colback=white,left=2pt,right=2pt,top=2pt,bottom=2pt,drop fuzzy shadow=black!25]\includegraphics[width=\linewidth]{#1}\end{tcolorbox}
   \vspace{-3pt}{\small\color{wdgray}#3}\end{figure}}
\newcommand{\plainfig}[3]{\par\noindent\begin{minipage}{\linewidth}\centering\includegraphics[width=#2\linewidth]{#1}\par\vspace{2pt}{\small\color{wdgray}#3}\end{minipage}\par\vspace{4pt}}

\begin{document}
% ---------- 封面条（仿课件封面） ----------
\thispagestyle{plain}
\begin{tcolorbox}[enhanced,enlarge left by=-4.8cm,width=\linewidth+4.8cm,arc=0pt,boxrule=0pt,
   colback=wddark,interior style={left color=wddark,right color=wdblue},left=18pt,right=18pt,top=14pt,bottom=12pt,
   overlay={\draw[white,line width=2.5pt] ([xshift=10pt,yshift=-10pt]frame.north west) -- ++(0,-0.9cm);
            \draw[white,line width=2.5pt] ([xshift=10pt,yshift=-10pt]frame.north west) -- ++(1.6cm,0);
            \draw[white,line width=2.5pt] ([xshift=-10pt,yshift=10pt]frame.south east) -- ++(0,0.9cm);
            \draw[white,line width=2.5pt] ([xshift=-10pt,yshift=10pt]frame.south east) -- ++(-1.6cm,0);}]
  {\color{white!85!wdlightblue}\small 数据结构 · 第3章
栈、队列和数组 · 康奈尔听课笔记}\par\vspace{6pt}
  {\color{white}\fontsize{22}{28}\selectfont\bfseries 3.2
课后习题：队列}\par\vspace{6pt}
  {\color{white!80!wdlightblue}\small 王道书 §3.2.5 试题精选 · §3.2.6
答案与解析（第 82--87 页）· 习题讲解字幕 P25--P29（选择 1--24）}
\end{tcolorbox}
\vspace{4pt}
\begin{tcolorbox}[enhanced,arc=3pt,boxrule=0.8pt,colframe=wdblue,colback=wdblue!3,left=10pt,right=10pt,top=6pt,bottom=6pt,
   colbacktitle=wdblue,coltitle=white,fonttitle=\bfseries\small,title={目录 · 点击标题跳转}]
\hypertarget{toc}{}\makeatletter\@starttoc{toc}\makeatother
\end{tcolorbox}
\hypertarget{ux8fd9ux7ec4ux9898ux5728ux8003ux4ec0ux4e48}{%
\section{这组题在考什么}\label{ux8fd9ux7ec4ux9898ux5728ux8003ux4ec0ux4e48}}

\tocback

\Needspace{4\baselineskip}

\cuenote{24 道题分成哪几类？循环队列的题，习题讲解老师让怎么做？}24
道选择题可以分成四类：

\begin{itemize}
\tightlist
\item
  \textbf{概念}（01--04）：队列和栈的区别、先进先出、允许的操作；
\item
  \textbf{循环队列}（05--10、21、22）：入队语句、队列长度、指针移动、判空判满。指针指向哪里每题都不一样，老师反复强调\hlred{不要记公式，画一个元素的情况倒推}；
\item
  \textbf{链式队列}（11--17）：哪种链表适合做队列、入队出队要改哪些指针；
\item
  \textbf{出队序列}（18--20、23、24）：两个队列、受限的双端队列、队列加栈，判断哪个序列得不到。
\end{itemize}

{\rowcolors{2}{wdblue!5}{white}\arrayrulecolor{wdblue!35}
\begin{xltabular}{\linewidth}{>{\hsize=0.474\hsize\raggedright\arraybackslash}X>{\hsize=2.053\hsize\raggedright\arraybackslash}X>{\hsize=0.474\hsize\raggedright\arraybackslash}X}
\rowcolor{wdblue}\textcolor{white}{\bfseries 题号} & \textcolor{white}{\bfseries 考点} & \textcolor{white}{\bfseries 答案} \\
\endfirsthead
01 & 栈和队列的主要区别 & D \\ \hline
02 & “先进先出”指什么 & B \\ \hline
03 & 允许对队列进行的操作 & D \\ \hline
04 & 入队 1234 的出队顺序 & B \\ \hline
05 & \texttt{A{[}0..n{]}} 循环队列的入队语句 & D \\ \hline
06 & \texttt{A{[}21{]}}，\texttt{front=8}、\texttt{rear=3} 时的长度 & C \\ \hline
07 & \texttt{A{[}0..5{]}}，删 1 个加 2 个后的指针 & B \\ \hline
08 & \texttt{front} 在队首前一个、\texttt{rear} 在队尾，判空 & D \\ \hline
09 & 没有其他数据成员时的判满条件 & C \\ \hline
10 & 用 \texttt{(rear+1)\%(n+1)==front} 判满，则 & A \\ \hline
11 & 链式队列和顺序队列相比 & D \\ \hline
12 & 最适合做队列的链表 & B \\ \hline
13 & 最不适合做链式队列的链表 & A \\ \hline
14 & 单链表实现队列，队头设在哪 & A \\ \hline
15 & 链式队列删除时要改哪个指针 & D \\ \hline
16 & 链式队列入队的语句 & D \\ \hline
17 & 队头在表尾、只设头指针的循环单链表，入队复杂度 & A \\ \hline
18 & 两个队列，不可能的输出序列 & B \\ \hline
19 & 输入受限、输出受限都得不到的序列 & C \\ \hline
20 & 【2010 统考】输出受限的双端队列 & C \\ \hline
21 & 【2011 统考】\texttt{front}、\texttt{rear} 的初值 & B \\ \hline
22 & 【2014 统考】\texttt{end1}、\texttt{end2} 的判空判满 & A \\ \hline
23 & 【2018 统考】队列加栈，不能得到的输出序列 & C \\ \hline
24 & 【2021 统考】输出受限的双端队列 & D \\ \hline
\end{xltabular}\par\vspace{4pt}}

\begin{fixbox}

\textbf{习题讲解视频的观看建议}：

\begin{itemize}
\tightlist
\item
  \textbf{P26（06--10）最值得看}：老师在这里讲透了循环队列题的做法：不记公式，画出“只有一个元素”的情况，再倒推队空、队满。后面
  21、22 都用这个方法；
\item
  \textbf{P28 的 19、20}：受限双端队列的“相邻”小秘诀，以及 20
  题怎样用模拟证明 C 不可能；
\item
  \textbf{P29（21--24）}：21、22 是 P26 方法的应用，24
  又用了一次“相邻”秘诀，比较短；
\item
  \textbf{P27 的 12}：A、B 两种链表都是
  \(O(1)\)，老师逐句写出入队、出队语句比较，想看可以看；
\item
  \textbf{P25（01--05）}：概念题，看书上解析即可。
\end{itemize}

下面每题的视频框都记了老师的做法。

\end{fixbox}

\hypertarget{ux4e00ux961fux5217ux7684ux6982ux5ff50104}{%
\section{一、队列的概念（01--04）}\label{ux4e00ux961fux5217ux7684ux6982ux5ff50104}}

\tocback

\Needspace{8\baselineskip}\vspace{22pt}

\cueblock{\cuenote}{\relax{栈和队列都是操作受限的线性表，区别在哪？}}

\begin{questionbox}[01]

栈和队列的主要区别在于（　）。

A. 它们的逻辑结构不一样　　　B. 它们的存储结构不一样\\
C. 所包含的元素不一样　　　　D. 插入、删除操作的限定不一样

\end{questionbox}

\hlred{\textbf{答案：D}}

\begin{videobox}

从“栈和队列都是操作受限的线性表”这句话找答案：两者逻辑结构都是线性表，A
错；栈和队列都既可以顺序存储，也可以链式存储，存储结构要看具体应用场景，B
错；元素类型不同不是栈和队列的区别，两个栈的元素类型也可以不同，C
错。区别在“受限”的方式：栈只在一端插入、删除，队列一端入队、另一端出队。

\end{videobox}

\begin{jiexibox}

王道书第 84--85
页：栈和队列的逻辑结构都是线性结构，都可以采用顺序存储或链式存储，选项 A
和 B 错误，选项 C 显然也错误。只有选项 D
才是栈和队列的本质区别，限定表中插入和删除操作位置的不同。

\end{jiexibox}

\begin{recallbox}

3.1 课后习题 01：栈和队列具有相同的\uln{逻辑结构}。本题问的是不同点。

\end{recallbox}

\begin{questionbox}[02]

队列的“先进先出”特性是指（　）。

Ⅰ. 最后插入队列中的元素总是最后被删除\\
Ⅱ. 当同时进行插入、删除操作时，总是插入操作优先\\
Ⅲ. 每当有删除操作时，总要先做一次插入操作\\
Ⅳ. 每次从队列中删除的总是最早插入的元素

A. Ⅰ　　　　B. Ⅰ 和 Ⅳ　　　　C. Ⅱ 和 Ⅲ　　　　D. Ⅳ

\end{questionbox}

\hlred{\textbf{答案：B}}

\begin{videobox}

先进先出就像食堂排队：先排的先打到饭，最后排的最后打到饭，所以 Ⅰ、Ⅳ
对。Ⅱ
错：用户“同时”发出插入和删除请求，计算机收到它们总有先后，还是先来先服务，不是插入优先。Ⅲ
错：可以连续出队好几次，不需要每次先插入。

\end{videobox}

\begin{jiexibox}

王道书第 85
页：队列“先进先出”的特性表现在：先入队列的元素先出队列，后入队列的元素后出队列，入队对应的是插入操作，出队对应的是删除操作。说法
Ⅰ 和 Ⅳ 均正确。

\end{jiexibox}

\begin{questionbox}[03]

允许对队列进行的操作有（　）。

A. 对队列中的元素排序　　　　B. 取出最近入队的元素\\
C. 在队列元素之间插入元素　　D. 删除队首元素

\end{questionbox}

\hlred{\textbf{答案：D}}

\begin{videobox}

A
和队列无关。B：最近入队的是队尾元素，如果先出它就成了后进先出。C：队列只在两端操作，不能在中间插入。D
就是出队。

\end{videobox}

\begin{jiexibox}

王道书第 85 页：删除队首元素即出队，是队列的基本操作之一，所以选择选项
D。

\end{jiexibox}

\begin{recallbox}

正课
3.2.1“三、队列的基本操作”：\uln{不允许直接访问或修改队列中间的元素}（王道书第
76 页）。

\end{recallbox}

\begin{questionbox}[04]

一个队列的入队顺序是 1, 2, 3, 4，则出队的输出顺序是（　）。

A. 4, 3, 2, 1　　　B. 1, 2, 3, 4　　　C. 1, 4, 3, 2　　　D. 3, 2, 4, 1

\end{questionbox}

\hlred{\textbf{答案：B}}

\begin{videobox}

先进先出，入队 1234，出队一定是 1234。即使入队、出队交替进行（1 进、1
出、2 进、3 进\ldots\ldots），出队序列还是
1234。这和栈不同：一个入栈序列可以有多种出栈序列，一个入队序列只有一种出队序列。

\end{videobox}

\begin{jiexibox}

王道书第 85 页：队列的入队顺序和出队顺序是一致的，这是和栈不同的。

\end{jiexibox}

\hypertarget{ux4e8cux5faaux73afux961fux521705102122}{%
\section{二、循环队列（05--10、21、22）}\label{ux4e8cux5faaux73afux961fux521705102122}}

\tocback

\Needspace{4\baselineskip}

\cuenote{看到一道循环队列的判空、判满、求长度题，第一步画什么？}习题讲解老师在
P26、P29 反复说的做法（原话“不要去记公式”）：

\begin{enumerate}
\def\labelenumi{\arabic{enumi}.}
\tightlist
\item
  先看数组下标范围，确定容量：\texttt{A{[}21{]}} 容量 21、下标
  0～20；\texttt{A{[}0..n{]}} 容量 \(n+1\)；
\item
  记住两条移动规则：\uln{入队一定移动 \texttt{rear}，出队一定移动
  \texttt{front}，都往下标增大的方向走}，走到末尾再回到 0（循环）；
\item
  求判空条件：\hlred{先画“队列里只有一个元素”的情况}（比如放在下标
  0），按题目说的指向画出
  \texttt{front}、\texttt{rear}；再倒推一步：这个元素进队前
  \texttt{rear} 在哪，那就是空队列的样子；
\item
  求判满条件：从空队列让元素一个个进队，进到题目说的最多个数为止，看两个指针的关系；
\item
  取模对数组长度（容量）取。
\end{enumerate}

理由是：题目里 \texttt{front}、\texttt{rear}
指向哪里每次都可能不一样，还可能加额外条件，公式记不住也用不上，要在考场上自己推。

\begin{recallbox}

正课 3.2.2 给的是公式和设计：

\begin{itemize}
\tightlist
\item
  方案一（牺牲一个单元，\texttt{front} 指队头、\texttt{rear}
  指队尾的下一个位置）：队空 \texttt{Q.rear==Q.front}，队满
  \texttt{(Q.rear+1)\%MaxSize==Q.front}，长度
  \texttt{(rear+MaxSize-front)\%MaxSize}；
\item
  六“rear 指向队尾元素”：初始化 \texttt{front=0,\ rear=MaxSize-1}，判空
  \texttt{(rear+1)\%MaxSize==front}，判满
  \texttt{(rear+2)\%MaxSize==front}。
\end{itemize}

正课是先讲原理再给结论；习题讲解老师主张做题时不背这些，按上面的步骤现场推。两种都记下，做题时用哪种都行。

\end{recallbox}

\begin{questionbox}[05]

循环队列存储在数组 \texttt{A{[}0..n{]}} 中，入队时的操作为（　）。

A. \texttt{rear=rear+1}　　　　　　　　B.
\texttt{rear=(rear+1)\ mod\ (n-1)}\\
C. \texttt{rear=(rear+1)\ mod\ n}　　　　D.
\texttt{rear=(rear+1)\ mod\ (n+1)}

\end{questionbox}

\hlred{\textbf{答案：D}}

\begin{videobox}

\texttt{A{[}0..n{]}} 下标 0～\(n\)，容量 \(n+1\)。入队移动队尾指针
\texttt{rear+1}；为了不越界，\texttt{rear} 在 \(n\) 时加 1 要回到
0，所以对容量 \(n+1\)
取模。老师举例：\(n=9\)，\texttt{rear=9}，\((9+1)\bmod 10=0\)；像只能保留一位数字的十进制加法，9+1
进位的 1 被舍掉，结果是 0。

\end{videobox}

\begin{jiexibox}

王道书第 85 页：数组下标范围为 \(0\sim n\)，因此数组容量为
\(n+1\)。循环队列中元素入队的操作是
\texttt{rear=(rear+1)\ mod\ maxsize}，题中
\texttt{maxsize=n+1}。因此入队操作应为
\texttt{rear=(rear+1)\ mod\ (n+1)}。

\end{jiexibox}

\Needspace{8\baselineskip}\vspace{22pt}

\cueblock{\cuenote}{\relax{不用公式，怎么数出 06 题的长度？}}

\begin{questionbox}[06]

已知循环队列的存储空间为数组 \texttt{A{[}21{]}}，\texttt{front}
指向队首元素的前一个位置，\texttt{rear} 指向队尾元素，假设当前
\texttt{front} 和 \texttt{rear} 的值分别为 8 和
3，则该队列的长度为（　）。

A. 5　　　　B. 6　　　　C. 16　　　　D. 17

\end{questionbox}

\hlred{\textbf{答案：C}}

\begin{videobox}

\texttt{A{[}21{]}} 容量 21、下标 0～20。从 \texttt{front}
往下标增大的方向一直到 \texttt{rear}，这一段放的是元素（\texttt{front}
本身不放，\texttt{rear} 放）。\texttt{rear=3}、\texttt{front=8}，下标
4～8 这 5 个位置是空的，所以长度 \(21-5=16\)。

想要公式，就现场推：找一个简单情况，假设 \texttt{rear=9}，此时只有下标 9
一个元素，长度 \(9-8=1\)，即 \texttt{rear-front}；\texttt{rear} 比
\texttt{front} 小时会出负数，就加上容量 21 再对 21
取模：\texttt{(rear-front+21)\%21}。老师说要掌握的是推公式的过程，不要像书上那样从原理一步步推，考场上来不及。

\end{videobox}

\slidefig{assets/3.2ex-06v.jpg}{0.6}{习题讲解视频 P26 板书：front、rear
一定往下标增大的方向移动；数出空位，21−5=16；再用 rear=9
的简单情况推出公式}

\begin{jiexibox}

王道书第 85 页：队列的长度为
\texttt{(rear-front+maxsize)\%maxsize\ =\ (rear-front+21)\%21\ =\ 16}。这种情况和
\texttt{front} 指向当前元素，\texttt{rear}
指向队尾元素的下一个元素的计算是相同的。

\textbf{注意}：数组 \texttt{A{[}n{]}} 的下标范围为 \(0\sim n-1\)。若写成
\texttt{A{[}0..n{]}}，则说明下标范围为 \(0\sim n\)。

\end{jiexibox}

\begin{questionbox}[07]

若用数组 \texttt{A{[}0..5{]}} 实现循环队列，且当前 \texttt{rear} 和
\texttt{front} 的值分别为 1 和
5，当从队列中删除一个元素，再加入两个元素后，\texttt{rear} 和
\texttt{front} 的值分别为（　）。

A. 3 和 4　　　　B. 3 和 0　　　　C. 5 和 0　　　　D. 5 和 1

\end{questionbox}

\hlred{\textbf{答案：B}}

\begin{videobox}

题目没说 \texttt{front}、\texttt{rear}
指向哪里，所以算不出队里现在有几个元素，但不影响做题：出队一定是
\texttt{front} 往后走，入队一定是 \texttt{rear} 往后走。容量
6，删除一个：\texttt{front} 从 5 到 0；加入两个：\texttt{rear} 从 1 到
3。

\end{videobox}

\slidefig{assets/3.2ex-07v.jpg}{0.6}{习题讲解视频 P26 板书：删除一个后 front 到 0，加入两个后 rear 到 3}

\begin{jiexibox}

王道书第 85 页：循环队列中，每删除一个元素，队首指针
\texttt{front=(front+1)\%6}，每插入一个元素，队尾指针
\texttt{rear=(rear+1)\%6}。上述操作后，\texttt{front=0}，\texttt{rear=3}。

\end{jiexibox}

\Needspace{8\baselineskip}\vspace{22pt}

\cueblock{\cuenote}{\relax{“一个元素倒推”在 08 题怎么用？}}

\begin{questionbox}[08]

假设用数组 \texttt{Q{[}MaxSize{]}} 实现循环队列，队首指针 \texttt{front}
指向队首元素的前一个位置，队尾指针 \texttt{rear}
指向队尾元素，则判断该队列为空的条件是（　）。

A. \texttt{Q.rear\ ==\ (Q.front+1)\%MaxSize}　　　B.
\texttt{(Q.rear+1)\%MaxSize\ ==\ Q.front+1}\\
C. \texttt{(Q.rear+1)\%MaxSize\ ==\ Q.front}　　　D.
\texttt{Q.rear\ ==\ Q.front}

\end{questionbox}

\hlred{\textbf{答案：D}}

\begin{videobox}

老师称这种做法为“逆向工程”：假设只有一个元素，在下标 0。\texttt{front}
指向它的前一个位置，下标 0 往前是 \texttt{MaxSize-1}；\texttt{rear}
指向这个元素，在 0。从空队列到这个状态，只做了一次入队：\texttt{front}
不动，\texttt{rear} 往后走了一步。所以空队列时 \texttt{rear} 在
\texttt{MaxSize-1}，和 \texttt{front} 在同一个地方，判空
\texttt{front==rear}。

老师也提到正课学过的三种区分队空队满的方法（牺牲一个单元、\texttt{size}、\texttt{tag}），说学的时候理解原理，做题时不要套里面的公式和结论。

\end{videobox}

\slidefig{assets/3.2ex-08av.jpg}{0.6}{习题讲解视频 P26 板书：只有一个元素（下标 0）时，front 在
MaxSize−1，rear 在 0}

\slidefig{assets/3.2ex-08bv.jpg}{0.6}{习题讲解视频 P26 板书：倒推出空队列，front 和 rear 指向同一处}

\begin{jiexibox}

王道书第 85 页：当队列中只有一个元素时，\texttt{front}
指向该元素的前一个位置，\texttt{rear}
指向该元素，因此当队列为空时，队首指针等于队尾指针，这样第一个元素入队后，才能符合题目要求。

\end{jiexibox}

\begin{questionbox}[09]

假设循环队列 \texttt{Q{[}MaxSize{]}} 的队首指针为
\texttt{front}，队尾指针为 \texttt{rear}，队列的最大容量为
\texttt{MaxSize}，此外，该队列再没有其他数据成员，则判断该队列已满的条件是（　）。

A. \texttt{Q.front\ ==\ Q.rear}　　　　　　　B.
\texttt{Q.front+Q.rear\ \textgreater{}=\ MaxSize}\\
C. \texttt{Q.front\ ==\ (Q.rear+1)\%MaxSize}　D.
\texttt{Q.rear\ ==\ (Q.front+1)\%MaxSize}

\end{questionbox}

\hlred{\textbf{答案：C}}

\begin{videobox}

“再没有其他数据成员”排除了 \texttt{size}、\texttt{tag}
两种方法，只能是牺牲一个存储单元。题目说的“最大容量
\texttt{MaxSize}”只能理解为牺牲之前数组的容量，老师说这题表达不太清楚：也没说
\texttt{front}、\texttt{rear}
指向哪里，是默认按书上的设计（\texttt{front} 指队头元素，\texttt{rear}
指队尾元素的后一个位置）。按这个设计画出队满：只空一个位置，\texttt{rear}
再加 1 就到 \texttt{front}，即
\texttt{(Q.rear+1)\%MaxSize==Q.front}。为什么那个位置不能用：再入队一个，\texttt{rear}
就和 \texttt{front} 重合，和队空的条件一样了。老师说 408
不会这么出，简单理解即可。

\end{videobox}

\slidefig{assets/3.2ex-09av.jpg}{0.6}{习题讲解视频 P26 板书：牺牲一个单元的队满状态，front 在 0，rear 在
MaxSize−1}

\slidefig{assets/3.2ex-09bv.jpg}{0.6}{习题讲解视频 P26 板书：如果最后一个单元也放元素，rear 回到 0 和 front
重合，就和队空分不开了}

\begin{fixbox}

习题讲解字幕 P26
里老师读到的题干是“已满足的条件”，老师说“足”字应改成“的”。看王道书第 82
页原题，已经印作“已满的条件”，题干没有问题。

\end{fixbox}

\begin{jiexibox}

王道书第 85
页：既然不能附加任何其他数据成员，只能采用牺牲一个存储单元的方法来区分是队空还是队满，约定以“队列头指针在队尾指针的下一位置作为队满的标志”，因此选择选项
C。选项 A 是判断队列是否空的条件，选项 B 和 D 都是干扰项。

\textbf{注意}：对于这类具体问题，举一些特例判断往往比直接思考问题能更快得到答案。

\end{jiexibox}

\begin{recallbox}

正课 3.2.2“五、方案一：牺牲一个存储单元”：队满
\texttt{(Q.rear+1)\%MaxSize==Q.front}，队空 \texttt{Q.rear==Q.front}。

\end{recallbox}

\begin{questionbox}[10]

假设用 \texttt{A{[}0..n{]}} 实现循环队列，\texttt{front}、\texttt{rear}
分别指向队首元素的前一个位置和队尾元素。若用
\texttt{(rear+1)\%(n+1)==front} 作为队满标志，则（　）。

A. 可用 \texttt{front==rear} 作为队空标志\\
B. 队列中最多可有 \(n+1\) 个元素\\
C. 可用 \texttt{front\textgreater{}rear} 作为队空标志\\
D. 可用 \texttt{(front+1)\%(n+1)==rear} 作为队空标志

\end{questionbox}

\hlred{\textbf{答案：A}}

\begin{videobox}

容量 \(n+1\)。先画一个元素在下标 0：\texttt{front} 在
\(n\)，\texttt{rear} 在 0。不断入队，\texttt{rear} 走到 \(n-1\) 时
\texttt{(rear+1)\%(n+1)==front}，只空着下标
\(n\)，说明用的是牺牲一个单元的方法。再倒推空队列：入队前 \texttt{rear}
在 \(n\)，和 \texttt{front} 重合，A 对。最多只能放 \((n+1)-1=n\)
个元素，B 错。A 对，C、D 自然不对。

\end{videobox}

\slidefig{assets/3.2ex-10av.jpg}{0.6}{习题讲解视频 P26 板书：一个元素在下标 0 时，front 在 n，rear 在 0}

\slidefig{assets/3.2ex-10bv.jpg}{0.6}{习题讲解视频 P26 板书：倒推出空队列，rear 和 front 都在 n}

\begin{jiexibox}

王道书第 85 页：若用 \texttt{(rear+1)\%(n+1)==front}
作为队满标志，则说明题目采用了牺牲一个存储单元的方法来区分队空和队满，因此可用
\texttt{front==rear} 作为队空标志。

\end{jiexibox}

\begin{questionbox}[21]

【2011 统考真题】已知循环队列存储在一维数组 \texttt{A{[}0..n-1{]}}
中，且队列非空时 \texttt{front} 和 \texttt{rear}
分别指向队首元素和队尾元素。若初始时队列为空，且要求第一个进入队列的元素存储在
\texttt{A{[}0{]}} 处，则初始时 \texttt{front} 和 \texttt{rear}
的值分别是（　）。

A. 0, 0　　　B. 0, \(n-1\)　　　C. \(n-1\), 0　　　D. \(n-1\), \(n-1\)

\end{questionbox}

\hlred{\textbf{答案：B}}

\begin{videobox}

老师又强调一次不要记公式，题目问法都很灵活。问初始值，就是问队空时的样子，所以先画第一个元素进队后的情况：元素在
\texttt{A{[}0{]}}，它既是队头也是队尾，\texttt{front}、\texttt{rear}
都指向 0。再倒推：入队会让 \texttt{rear}
往下标增大的方向走一步，所以进队前 \texttt{rear} 在 0
的前一个位置，循环回去就是 \(n-1\)；\texttt{front} 不动，还是 0。

\end{videobox}

\slidefig{assets/3.2ex-21v.jpg}{0.6}{习题讲解视频 P29 板书：第一个元素进队后 front、rear 都在 0，倒推出 rear
初值为 n−1}

\begin{jiexibox}

王道书第 86 页：根据题意，第一个元素进入队列后存储在 \texttt{A{[}0{]}}
处，此时 \texttt{front} 和 \texttt{rear} 值都为 0。入队时因为要执行
\texttt{(rear+1)\%n} 操作，所以若入队后指针指向 0，则 \texttt{rear}
初值为 \(n-1\)，而因为第一个元素在 \texttt{A{[}0{]}} 中，插入操作只改变
\texttt{rear} 指针，所以 \texttt{front} 为 0 不变。

\textbf{注意}：①
循环队列是指顺序存储的队列，而不是指逻辑上的循环，如循环单链表表示的队列不能称为循环队列。②
\texttt{front} 和 \texttt{rear} 的初值并不是固定的。

\end{jiexibox}

\begin{recallbox}

正课 3.2.2“六、rear 指向队尾元素”就是这种设计：初始化
\texttt{front=0,\ rear=MaxSize-1}，入队先移 \texttt{rear} 再放元素。

\end{recallbox}

\begin{questionbox}[22]

【2014 统考真题】循环队列放在一维数组 \texttt{A{[}0..M-1{]}}
中，\texttt{end1} 指向队首元素，\texttt{end2}
指向队尾元素的后一个位置。假设队列两端均可进行入队和出队操作，队列中最多能容纳
\(M-1\) 个元素。初始时为空。下列判断队空和队满的条件中，正确的是（　）。

{\rowcolors{2}{wdblue!5}{white}\arrayrulecolor{wdblue!35}
\begin{xltabular}{\linewidth}{>{\hsize=0.162\hsize\raggedright\arraybackslash}X>{\hsize=1.297\hsize\raggedright\arraybackslash}X>{\hsize=1.541\hsize\raggedright\arraybackslash}X}
\rowcolor{wdblue}\textcolor{white}{\bfseries } & \textcolor{white}{\bfseries 队空} & \textcolor{white}{\bfseries 队满} \\
\endfirsthead
A. & \texttt{end1\ ==\ end2} & \texttt{end1\ ==\ (end2+1)\ mod\ M} \\ \hline
B. & \texttt{end1\ ==\ end2} & \texttt{end2\ ==\ (end1+1)\ mod\ (M-1)} \\ \hline
C. & \texttt{end2\ ==\ (end1+1)\ mod\ M} & \texttt{end1\ ==\ (end2+1)\ mod\ M} \\ \hline
D. & \texttt{end1\ ==\ (end2+1)\ mod\ M} & \texttt{end2\ ==\ (end1+1)\ mod\ (M-1)} \\ \hline
\end{xltabular}\par\vspace{4pt}}

\end{questionbox}

\hlred{\textbf{答案：A}}

\begin{videobox}

同样的步骤。一个元素在下标 0：\texttt{end1} 指向它，在 0；\texttt{end2}
指向它的后一个位置，在 1。入队让 \texttt{end2} 走了一步，所以空队列时
\texttt{end2} 在 0，和 \texttt{end1} 相同，队空 \texttt{end1==end2}，在
A、B 中选。队满：从空队列连续入队 \(M-1\) 个，\texttt{end2} 走到
\(M-1\)，此时 \texttt{(end2+1)\ mod\ M\ ==\ end1}。取模对数组长度 \(M\)
取，不是 \(M-1\)，选 A。

\end{videobox}

\slidefig{assets/3.2ex-22av.jpg}{0.6}{习题讲解视频 P29 板书：由一个元素的情况倒推出空队列，end1、end2 都在 0}

\slidefig{assets/3.2ex-22bv.jpg}{0.6}{习题讲解视频 P29 板书：入队 M−1 个后 end2 在 M−1，队满 end1==(end2+1)
mod M}

\begin{jiexibox}

王道书第 87 页：\texttt{end1} 指向队首元素，可知出队操作是先从
\texttt{A{[}end1{]}} 读数，然后 \texttt{end1} 再加 1。\texttt{end2}
指向队尾元素的后一个位置，可知入队操作是先存数到
\texttt{A{[}end2{]}}，然后 \texttt{end2} 再加 1。若用 \texttt{A{[}0{]}}
存储第一个元素，则队列初始时，入队操作先把数据放到 \texttt{A{[}0{]}}
中，然后 \texttt{end2} 自增，即可知 \texttt{end2} 初值为 0；而
\texttt{end1} 指向的是队首元素，队首元素在数组 A 中的下标为 0，所以得知
\texttt{end1} 的初值也为 0，可知队空条件为
\texttt{end1==end2}；然后考虑队列满时，因为队列最多能容纳 \(M-1\)
个元素，假设队列存储在下标为 0 到 \(M-2\) 的 \(M-1\) 个区域，队头为
\texttt{A{[}0{]}}，队尾为
\texttt{A{[}M-2{]}}，此时队列满，考虑在这种情况下 \texttt{end1} 和
\texttt{end2} 的状态，\texttt{end1} 指向队首元素，可知
\texttt{end1=0}，\texttt{end2} 指向队尾元素的后一个位置，可知
\texttt{end2=M-2+1=M-1}，所以队满条件为
\texttt{end1==(end2+1)\ mod\ M}。

\end{jiexibox}

\begin{recallbox}

\texttt{end1}、\texttt{end2} 就是正课 3.2.2 方案一的
\texttt{front}、\texttt{rear}，结论和方案一相同。

\end{recallbox}

\Needspace{18\baselineskip}

\hypertarget{ux4e09ux94feux5f0fux961fux52171117}{%
\section{三、链式队列（11--17）}\label{ux4e09ux94feux5f0fux961fux52171117}}

\tocback

\Needspace{4\baselineskip}

\cuenote{判断一种链表适不适合做队列，看哪两个操作？}老师的做法：队列只有入队（在队尾）和出队（在队头），所以看这种链表能不能
\(O(1)\) 找到队头和队尾，再看入队、出队各要写几句。

\begin{questionbox}[11]

与顺序队列相比，链式队列的（　）。

A. 优点是队列的长度不受限制\\
B. 优点是入队和出队时间效率更高\\
C. 缺点是不能进行顺序访问\\
D. 缺点是不能根据队首指针和队尾指针计算队列的长度

\end{questionbox}

\hlred{\textbf{答案：D}}

\begin{videobox}

A 太绝对：链式队列可以随时申请新结点，但还受内存空间限制。B
错：两者都是改指针或数组下标，入队、出队效率都很高。C
错：顺序队列和链式队列都是队列，只能队尾进、队头出，\uln{都不能顺序访问}。D
对：链式队列的结点不连续存放，从队头、队尾两个结点的地址算不出中间有几个元素；顺序队列的
\texttt{front}、\texttt{rear} 是数组下标，可以算。

\end{videobox}

\begin{jiexibox}

王道书第 85
页：虽然链式队列采用动态分配方式，但其长度也受内存空间的限制，不能无限制增长。顺序队列和链式队列的入队和出队时间复杂度均为
\(O(1)\)。\uln{顺序队列和链式队列都可以进行顺序访问}。对于顺序队列，可通过队首指针和队尾指针计算队列中的元素个数，而链式队列则不能。

\end{jiexibox}

\begin{addbox}

C
选项老师和书的理由相反：老师说两者都\uln{不能}顺序访问，书上说两者都\uln{可以}顺序访问。结论一样：这不是链式队列独有的缺点，C
错。两种说法都按原样记下。

\end{addbox}

\Needspace{8\baselineskip}\vspace{22pt}

\cueblock{\cuenote}{\relax{A、B 都是 \(O(1)\)，为什么选 B？}}

\begin{questionbox}[12]

下列描述的几种链表中，最适合用作队列的是（　）。

A. 带队首指针和队尾指针的循环单链表\\
B. 带队首指针和队尾指针的非循环单链表\\
C. 只带队首指针的非循环单链表\\
D. 只带队首指针的循环单链表

\end{questionbox}

\hlred{\textbf{答案：B}}

\begin{videobox}

老师把 A、B
的入队、出队语句都写了出来（按不带头结点画，带不带都不影响）：

\begin{itemize}
\tightlist
\item
  A
  入队要三句：\texttt{rear-\textgreater{}next=q;\ q-\textgreater{}next=front;\ rear=q;}，多出的一句是为了保持循环；出队也要先改
  \texttt{rear-\textgreater{}next} 再移动 \texttt{front}；
\item
  B
  入队两句：\texttt{rear-\textgreater{}next=q;\ rear=q;}；出队两句：保存
  \texttt{front} 的数据，\texttt{front=front-\textgreater{}next}。
\end{itemize}

两者都是 \(O(1)\)，但 B 更简单，A“画蛇添足”“杀鸡用牛刀”。C、D
没有队尾指针，入队要先遍历找到队尾，\(O(n)\)，排除。

\end{videobox}

\begin{jiexibox}

王道书第 85--86 页：因为队列需在双端进行操作，所以选项 C 和 D
的链表显然不太适合链队。对于选项
A，链表在完成入队和出队后还要修改为循环的，对于队列来讲这是多余的（画蛇添足）。对于选项
B，因为有首指针，所以适合删除首结点；因为有队尾指针，所以适合在其后插入结点。

\end{jiexibox}

\begin{questionbox}[13]

下列描述的几种链表中，最不适合用作链式队列的是（　）。

A. 只带队首指针的非循环双链表　　B. 只带队首指针的循环双链表\\
C. 只带队尾指针的循环双链表　　　D. 只带队尾指针的循环单链表

\end{questionbox}

\hlred{\textbf{答案：A}}

\begin{videobox}

A：只有队首指针、又不循环，入队前要从队头遍历到队尾，\(O(n)\)。B、C：循环双链表功能最强，有队头或队尾任何一个指针，都能
\(O(1)\) 找到另一头，入队出队都是 \(O(1)\)。D：队尾指针 \texttt{p} 的
\texttt{next} 就是队头。入队
\texttt{q-\textgreater{}next=p-\textgreater{}next;\ p-\textgreater{}next=q;\ p=q;}，出队通过
\texttt{p-\textgreater{}next-\textgreater{}next} 改指针，也都是
\(O(1)\)。

\end{videobox}

\begin{jiexibox}

王道书第 86
页：因为非循环双链表只带队首指针，所以在执行入队操作时需要修改队尾结点的指针域，而查找队尾结点需要
\(O(n)\) 的时间。选项 B、C 和 D 均可在 \(O(1)\) 的时间内找到队首和队尾。

\end{jiexibox}

\begin{recallbox}

正课 2.3.4 循环链表：循环单链表\uln{仅设尾指针 \texttt{r}}
时，在表头或表尾插入都是
\(O(1)\)；若用头指针，表尾插入要遍历整个链表，\(O(n)\)。D 选项和 17
题都是这个道理。

\end{recallbox}

\begin{questionbox}[14]

在用单链表实现队列时，队头设在链表的（　）位置。

A. 链头　　　　B. 链尾　　　　C. 链中　　　　D. 以上都可以

\end{questionbox}

\hlred{\textbf{答案：A}}

\begin{videobox}

队头用来出队，从出队考虑：队头在链头，出队后指针往后移一个就是新队头，方便；队头在链尾，新队头是它的前驱，单链表找前驱要从头遍历，不行；设在链中，连哪里是队头、哪里是队尾都说不清，不可能这么设计。

\end{videobox}

\begin{jiexibox}

王道书第 86
页：因为在队头做出队操作，为便于删除队首元素，所以总是选择链头作为队头。

\end{jiexibox}

\begin{questionbox}[15]

用链式存储方式的队列进行删除操作时需要（　）。

A. 仅修改头指针　　　　　　　B. 仅修改尾指针\\
C. 头尾指针都要修改　　　　　D. 头尾指针可能都要修改

\end{questionbox}

\hlred{\textbf{答案：D}}

\begin{videobox}

没说是哪种链表，默认最普通的单链表，老师按不带头结点画。出队一定要把
\texttt{front} 往后移。特殊情况：队里只有一个元素时，\texttt{front} 和
\texttt{rear} 指向同一个结点，它出队后 \texttt{front}
指向空，\texttt{rear} 也要改成指向空，否则 \texttt{rear}
还指着已经删掉的结点。所以头尾指针可能都要改。

\end{videobox}

\begin{jiexibox}

王道书第 86
页：队列采用链式存储时，删除元素要从表头删除，通常仅需修改头指针，但若队列中仅有一个元素，则队尾指针也需要被修改，仅有一个元素时，删除后队列为空，需要修改队尾指针为
\texttt{rear=front}。

\end{jiexibox}

\begin{recallbox}

正课 3.2.3“四、出队”：带头结点时，最后一个结点出队要
\texttt{Q.rear=Q.front}（书上解析说的就是这种）；不带头结点时，要把
\texttt{front} 和 \texttt{rear} 都置为 \texttt{NULL}（老师画的是这种）。

\end{recallbox}

\begin{questionbox}[16]

在一个链式队列中，假设队首指针为 \texttt{front}，队尾指针为
\texttt{rear}，\texttt{x}
所指向的元素需要入队，则需要执行的操作为（　）。

A. \texttt{front=x,\ front=front-\textgreater{}next}　　　B.
\texttt{x-\textgreater{}next=front-\textgreater{}next,\ front=x}\\
C. \texttt{rear-\textgreater{}next=x,\ rear=x}　　　　　　D.
\texttt{rear-\textgreater{}next=x,\ x-\textgreater{}next=NULL,\ rear=x}

\end{questionbox}

\hlred{\textbf{答案：D}}

\begin{videobox}

入队在队尾：\texttt{rear-\textgreater{}next=x} 把 \texttt{x}
接到队尾后面；\texttt{x} 成了新的队尾，\texttt{x-\textgreater{}next}
要指向空；最后 \texttt{rear=x} 移动队尾指针。

\end{videobox}

\begin{jiexibox}

王道书第 86 页：插入操作时，先将结点 \texttt{x} 插入到链表尾部，再让
\texttt{rear} 指向这个结点 \texttt{x}。选项 C
的做法不够严密，因为是队尾，所以队尾 \texttt{x-\textgreater{}next}
必须置为空。

\end{jiexibox}

\begin{recallbox}

正课 3.2.3“三、入队”的代码里也有 \texttt{s-\textgreater{}next=NULL}
这一句。

\end{recallbox}

\begin{questionbox}[17]

假设循环单链表表示的队列长度为
\(n\)，队头固定在链表尾，若只设头指针，则入队操作的时间复杂度为（　）。

A. \(O(n)\)　　　B. \(O(1)\)　　　C. \(O(n^2)\)　　　D. \(O(n\log_2 n)\)

\end{questionbox}

\hlred{\textbf{答案：A}}

\begin{videobox}

队头在链表尾，那么队尾就在链表头（头指针 \texttt{head}
指的是队尾）。入队要把新结点 \texttt{p} 放到 \texttt{head}
前面成为新队尾，还要保持循环：链表尾（队头）的 \texttt{next} 要改成指向
\texttt{p}。只有头指针，找链表尾要遍历一圈，\(O(n)\)；找到后
\texttt{p-\textgreater{}next=head;\ front-\textgreater{}next=p;\ head=p;}
三句是 \(O(1)\)。合起来
\(O(n)\)。老师特别提醒：\uln{队头不一定是链表头，题目给的指针也不一定是队头指针}，要看题目说明，最好画图。

\end{videobox}

\begin{jiexibox}

王道书第 86
页：依题意，入队操作是在队尾进行，即链表表头。题中已明确说明链表只设头指针，即没有头结点和尾指针，入队后，循环单链表必须保持循环的性质，在只带头指针的循环单链表中寻找表尾结点的时间复杂度为
\(O(n)\)，所以入队的时间复杂度为 \(O(n)\)。

\end{jiexibox}

\begin{addbox}

王道书第 86 页 21
题后的注意：循环单链表表示的队列\uln{不能称为循环队列}，循环队列专指顺序存储的队列。本题是“用循环单链表表示的队列”。

\end{addbox}

\hypertarget{ux56dbux51faux961fux5e8fux521718202324}{%
\section{四、出队序列（18--20、23、24）}\label{ux56dbux51faux961fux5e8fux521718202324}}

\tocback

\Needspace{8\baselineskip}\vspace{22pt}

\cueblock{\cuenote}{\relax{受限双端队列的“相邻”秘诀是什么？它能证明一个序列“不可能”吗？}}

\begin{questionbox}[18]

假设输入序列为 1, 2, 3, 4,
5，利用两个队列进行出入队操作，不可能输出的序列是（　）。

A. 1, 2, 3, 4, 5　　B. 5, 2, 3, 4, 1　　C. 1, 3, 2, 4, 5　　D. 4, 1, 5,
2, 3

\end{questionbox}

\hlred{\textbf{答案：B}}

\begin{videobox}

逐个模拟，设两个队列为 Q1、Q2。B：5 最后输入却第一个输出，只能是 1234
都进了一个队列、5 单独进另一个队列，那么 5 之后一定是 1234，不可能是
2341。C：1 随便进一个；2、3 分进两个队列；4、5 随便进，按 13245
出队即可。D：123 进 Q1，4 进 Q2，5 也进 Q2，按 41523
出队即可。老师说这类题只能手动模拟。

\end{videobox}

\begin{jiexibox}

王道书第 86 页：此类题可对各选项进行模拟，假设队列为 \(Q_1\) 和
\(Q_2\)。对于选项 A，元素依次入队 \(Q_1\)，然后依次出队。对于选项 B，5
最先出队，只可能是 1, 2, 3, 4 入队 \(Q_1\)，5 入队 \(Q_2\)，然后 5
出队，只能得到 5, 1, 2, 3, 4。对于选项 C，1, 2, 5 入队 \(Q_1\)，3, 4
入队 \(Q_2\)，然后按要求出队。选项 D 的分析同选项 C。

\end{jiexibox}

\begin{questionbox}[19]

若以 1, 2, 3, 4
作为双端队列的输入序列，则既不能由输入受限的双端队列得到，又不能由输出受限的双端队列得到的输出序列是（　）。

A. 1, 2, 3, 4　　　B. 4, 1, 3, 2　　　C. 4, 2, 3, 1　　　D. 4, 2, 1, 3

\end{questionbox}

\hlred{\textbf{答案：C}}

\begin{videobox}

B、C、D 都是 4 先输出，所以 1234 一定已经全部输入。

\begin{itemize}
\tightlist
\item
  \textbf{输入受限}（Q1，一端输入、两端输出）：1234 依次进队后，4
  从一端出去，下一个只能是另一端的 1 或者这一端的 3，不可能是 2。所以
  C、D 都得不到；B（4132）可以：4、1、3、2 分别从两端出；
\item
  \textbf{输出受限}（Q2，两端输入、一端输出）：全部进队后只能从一端出，所以队里从出口往里的排列就是输出序列。C
  要求排成 4231：1 先进，2
  无论从哪边进\hlred{都和 1 相邻}，中间不可能隔着 3。所以 C
  也得不到。D（4213）可以：1 进，2 从一边进，3 从另一边进，4 从 2
  那边进，排成 4213。
\end{itemize}

所以 C
两种都得不到。这里老师讲了一个常用技巧：从第一个元素开始，每个新元素都只能接在已有排列的左边或右边，看能不能拼出想要的排列。

\end{videobox}

\begin{jiexibox}

王道书第 86
页：使用排除法。先看可由输入受限的双端队列产生的序列：设右端输入受限，1,
2, 3, 4 依次左入，则依次左出可得 4, 3, 2, 1，排除选项
A；左出、右出、左出、左出可得到 4, 1, 3, 2，排除选项
B；再看可由输出受限的双端队列产生的序列：设右端输出受限，1, 2, 3, 4
依次左入、左入、右入、左入，依次左出可得到 4, 2, 1, 3，排除选项 D。

\end{jiexibox}

\begin{recallbox}

正课 3.2.4 的结论正是这三个选项：输入 1234 时，输入受限非法
4213、4231，输出受限非法
4132、4231；\uln{4132 只有输入受限能得到，4213 只有输出受限能得到，4231 两种都不行}。

\end{recallbox}

\begin{questionbox}[20]

【2010
统考真题】某队列允许在其两端进行入队操作，但仅允许在一端进行出队操作。若元素
\(a, b, c, d, e\)
依次入此队列后再进行出队操作，则不可能得到的出队序列是（　）。

A. \(b, a, c, d, e\)　　B. \(d, b, a, c, e\)　　C.
\(d, b, c, a, e\)　　D. \(e, c, b, a, d\)

\end{questionbox}

\hlred{\textbf{答案：C}}

\begin{videobox}

这是输出受限的双端队列。老师的\textbf{小秘诀}：入队顺序固定是
\(abcde\)，两端都能进，所以 \(a\) 进队后，\(b\) 一定挨着 \(a\)；\(ab\)
进队后，\(c\) 要么在 \(a\) 外侧、要么在 \(b\)
外侧\ldots\ldots 熟练了不用画队列，直接看选项能不能从 \(a\)
开始“往两边接”拼出来，能拼出来，就让所有元素进队后再全部从出口出。A、B、D
都能拼出（如 D：\(a\)，\(b\) 接在 \(a\) 旁，\(c\) 接在 \(b\) 旁，\(d\)
接另一边，\(e\) 接 \(c\) 那边）。C 里 \(b\) 和 \(a\) 之间隔着
\(c\)，拼不出来。

老师接着说：题目没有要求所有元素都进队之后才能出队，所以“拼不出来”还要证明。证明：\(d\)
第一个出队，\(a, b, c\) 已在队中。\(d\) 之后紧接着是 \(b\)，所以 \(b\)
必须在 \(a\) 靠出口的一侧；\(c\) 不能在 \(b\) 外侧（否则 \(d\)
后面先出的是 \(c\)），只能在 \(a\) 另一侧。这时队里从出口往里是
\(b, a, c\)，\(b\) 出队后接着只能是 \(a\)，不可能是
\(c\)。老师说一般用秘诀就够了，只有几个选项都判断不出来时，才这样细致地模拟。

\end{videobox}

\slidefig{assets/3.2ex-20v.jpg}{0.6}{习题讲解视频 P28 板书：证明 C 不可能。d、b 出队后队里是 a、c（e
未进），下一个只能是 a}

\begin{jiexibox}

王道书第 86 页：本题的队列实际上是一个输出受限的双端队列，如图 3.11
所示。A 操作：\(a\) 左入（或右入）、\(b\) 左入、\(c\) 右入、\(d\)
右入、\(e\) 右入。B 操作：\(a\) 左入（或右入）、\(b\) 左入、\(c\)
右入、\(d\) 左入、\(e\) 右入。D 操作：\(a\) 左入（或右入）、\(b\)
左入、\(c\) 左入、\(d\) 右入、\(e\) 左入。C 操作：\(a\)
左入（或右入）、\(b\) 右入、因 \(d\) 未出，此时只能入队，\(c\)
怎么进都不可能在 \(b\) 和 \(a\) 之间。

【另解】初始时队列为空，第 1 个元素 \(a\) 左入（或右入）后，第 2 个元素
\(b\) 无论是左入还是右入都必与 \(a\) 相邻，而选项 C 中 \(a\) 与 \(b\)
不相邻，不合题意。

\end{jiexibox}

\begin{addbox}

“全部进队再一起出队、看能不能往两边接出来”\uln{只是充分条件}：能拼出来，序列一定合法；拼不出来，不一定不合法，因为还可以边进边出。我用程序穷举验证：输出受限的双端队列，\(n=4\)
时“全部进队再出”只能得到 8 种序列，而边进边出共能得到 22 种；\(n=5\)
时是 16 种对 90 种。例如 1324：1 进、1 出，2 进，3 从出口一侧进、3 出，2
出，4 进、4 出，合法；但“全部进队再出”拼不出来（2 和 1 中间隔着 3）。

所以秘诀拼不出来时，要像老师证明 C 那样再模拟一下。本节 19、20、24
题各选项我都用程序穷举所有入队、出队顺序核对过，答案和书一致。

\end{addbox}

\begin{questionbox}[23]

【2018 统考真题】现有队列 Q 与栈 S，初始时 Q 中的元素依次是 1, 2, 3, 4,
5, 6（1 在队头），S 为空。若仅允许下列 3 种操作：① 出队并输出出队元素；②
出队并将出队元素入栈；③
出栈并输出出栈元素，则不能得到的输出序列是（　）。

A. 1, 2, 5, 6, 4, 3　　B. 2, 3, 4, 5, 6, 1　　C. 3, 4, 5, 6, 1, 2　　D.
6, 5, 4, 3, 2, 1

\end{questionbox}

\hlred{\textbf{答案：C}}

\begin{videobox}

逐个模拟。A：1、2 直接输出，3、4 进栈，5、6 直接输出，最后出栈
4、3。B：1 进栈，2～6 直接输出，最后出栈 1。C：第一个输出 3，所以 1、2
要先进栈，3456 输出后，栈里出来的顺序是 2、1，只能得到 345621，得不到
345612。D：1～6 全部进栈再全部出栈。

\end{videobox}

\begin{jiexibox}

王道书第 87 页：选项 A 的操作顺序为①①②②①①③③。选项 B
的操作顺序为②①①①①①③。选项 D 的操作顺序为②②②②②②③③③③③③。对于选项
C：首先输出 3，说明 1 和 2 必须先依次入栈，而此后 2 肯定比 1
先输出，因此无法得到 1, 2 的输出顺序。

\end{jiexibox}

\begin{recallbox}

3.1
课后习题第六节“小结与速查”的提醒：这类题加了额外条件，要先把条件翻译过来再判断。本题的队列只是按
1～6 的顺序送出元素，每个元素要么直接输出，要么进栈。

\end{recallbox}

\begin{questionbox}[24]

【2021 统考真题】初始为空的队列 Q
的一端仅能进行入队操作，另外一端既能进行入队操作又能进行出队操作。若 Q
的入队序列是 1, 2, 3, 4, 5，则不能得到的出队序列是（　）。

A. 5, 4, 3, 1, 2　　B. 5, 3, 1, 2, 4　　C. 4, 2, 1, 3, 5　　D. 4, 1, 3,
2, 5

\end{questionbox}

\hlred{\textbf{答案：D}}

\begin{videobox}

一端只能入队、另一端能入能出，就是只在一端输出的输出受限双端队列，用 20
题的秘诀：所有元素进队后再从同一端出。A：1，2 接一边，3、4、5
接另一边，排成 54312。B：1，2 接 1 的右边，3 接 1 的左边，4 接 2
的右边，5 接 3 的左边，排成 53124。C：同理排成 42135。D：2 要么在 1
左边、要么在 1 右边，中间不可能隔着 3。

老师也按模拟证明了 D：4 第一个出队，1、2、3 已在队中；4 出队后接着是
1，所以 2 只能在 1 远离出口的一侧，3 也只能接在 2 外侧。4、1
出队后，队里从出口往里是 2、3，一定是 2 先于 3 出队，得不到
41325。老师说考场上 A、B、C 瞄一眼就知道可以，D
也不用这样证明，能选出来就行。

\end{videobox}

\slidefig{assets/3.2ex-24v.jpg}{0.6}{习题讲解视频 P29 板书：证明 D 不可能。4、1 出队后队里是 2、3（5
未进），2 一定先于 3 出队}

\begin{jiexibox}

王道书第 87 页：假设队列左端允许入队和出队，右端只能入队。对于选项
A，依次从右端入队 1, 2，再从左端入队 3, 4, 5。对于选项 B，从右端入队 1,
2，然后从左端入队 3，再从右端入队 4，最后从左端入队 5。对于选项
C，从左端入队 1, 2，然后从右端入队 3，再从左端入队 4，最后从右端入队
5。无法验证选项 D 的序列。

\end{jiexibox}

\slidefig{assets/3.2ex-24-book.png}{0.75}{王道书第 87 页：24 题 A、B、C 全部入队后的队列（左端出队）}

\begin{recallbox}

正课
3.2.4“输出受限的双端队列”：要按某个顺序输出，这些元素在队中就必须按这个顺序排好，看能不能从两边交替插入拼出来；4132
不合法正是因为“2 一定和 1 相邻”。

\end{recallbox}

\newpage

\hypertarget{ux4e94ux603bux7ed3}{%
\section{五、总结}\label{ux4e94ux603bux7ed3}}

\tocback

\begin{summarybox}

\textbf{循环队列题不背公式，现场推}（05--10、21、22）：

\begin{itemize}
\tightlist
\item
  先看容量：\texttt{A{[}n{]}} 下标 \(0\sim n-1\)，\texttt{A{[}0..n{]}}
  下标 \(0\sim n\)、容量 \(n+1\)；取模对容量取；
\item
  入队动 \texttt{rear}、出队动
  \texttt{front}，都往下标增大的方向走，到末尾回到
  0；指针指向哪里不知道也能算指针的移动（07）；
\item
  判空：画只有一个元素的情况，倒推一次入队前的样子（08、10、21、22）；判满：从空队列入队到最多个数（09、10、22）；
\item
  求长度：数出空着的位置，用容量减（06）；
\item
  只有 \texttt{front}、\texttt{rear}
  两个成员时，只能牺牲一个单元区分队空队满（09）。
\end{itemize}

\textbf{链式队列}：

\begin{itemize}
\tightlist
\item
  看能不能 \(O(1)\) 找到队头和队尾；都能时，比较哪种语句更少（12
  选非循环单链表，循环是多余的）；
\item
  单链表的队头设在链头（14）；入队
  \texttt{rear-\textgreater{}next=x;\ x-\textgreater{}next=NULL;\ rear=x;}（16）；
\item
  只剩一个元素时出队，头尾指针都要改（15）；
\item
  队头不一定是链表头，按题目画图（17）；
\item
  链式队列不能由头尾指针算长度（11）；循环单链表表示的队列不叫循环队列（21
  注意）。
\end{itemize}

\textbf{出队序列}：

\begin{itemize}
\tightlist
\item
  两个队列、队列加栈：逐个模拟（18、23）；
\item
  输出受限：每个新元素只能接在已有排列的两边，“\(b\) 一定和 \(a\)
  相邻”（19、20、24）；拼得出就合法，拼不出还要模拟证明；
\item
  输入受限：全部进队后，每次只能从两端取（19）。
\end{itemize}

\textbf{自测}：

\begin{enumerate}
\def\labelenumi{\arabic{enumi}.}
\tightlist
\item
  \texttt{A{[}0..n{]}} 做循环队列，\texttt{front}
  指队头元素、\texttt{rear} 指队尾元素，第一个元素放在
  \texttt{A{[}0{]}}，初值是多少？
\item
  08 题的设计下，队满条件是什么（牺牲一个单元）？
\item
  为什么带头尾指针的循环单链表不如非循环单链表适合做队列？
\item
  用“相邻”秘诀说明 \(abcde\) 进输出受限队列，为什么得不到 \(dbcae\)？
\item
  1324 为什么能由输出受限的双端队列得到，却不能用“全部进队再出”拼出来？
\end{enumerate}

\end{summarybox}
\end{document}
